%
% teil3.tex -- Beispiel-File für Teil 3
%
% (c) 2020 Prof Dr Andreas Müller, Hochschule Rapperswil
%
% !TEX root = ../../buch.tex
% !TEX encoding = UTF-8
%
\section{Analytische Fortsetzung
\label{Analytische Fortsetzung}}
\kopfrechts{Analytische Fortsetzung}

Die analytische Fortsetzung ist ein Konzept aus der komplexen Analysis, mit dem eine holomorphe Funktion über ihren ursprünglichen Definitionsbereich hinaus erweitert werden kann.
Eine ausführliche Einführung dazu findet sich in Abschnitt 5.3 des Buches.
Grundsätzlich wird dabei eine neue Funktion konstruiert, die auf dem ursprünglichen Definitionsbereich mit der alten Funktion übereinstimmt, aber auch ausserhalb davon definiert ist.

In der Anwendung auf die riemannsche Zeta-Funktion kann eine solche analytische Fortsetzung fast auf die gesamte komplexe Ebene angewendet werden.
\index{analytische Fortsetzung}%
Die einzige Ausnahme ist der Punkt bei $s = 1$.
%Beschreiben anhand Bild wie so eine Funktion erweitert wird anhand prinzip dass man definitionsbereich von Funktionen sucht, die in einem gewissen Bereich sich überschneiden

\subsection{Anwendung auf die Zeta-Funktion\label{Anwendung auf die Zeta-Funktion}}

Wie bereits in Abschnitt~\ref{Riemannsche Zeta-Funktion} erwähnt, ist die Reihendarstellung
\index{zeta-Funktion@$\zeta$-Funktion!Reihendarstellung}%
\[
\zeta(s) = 1 + \frac{1}{2^{s}} + \frac{1}{3^{s}} + \frac{1}{4^{s}} + \ldots
\]
nur für $\operatorname{Re}(s) > 1$ korrekt.

Wie in Kapitel \ref{chapter:hankel} hergeleitet wird, liefert die analytische Fortsetzung für die Zeta-funktion die sogenannte Funktionalgleichung, welche $\zeta(s)$ mit $\zeta(1-s)$ verbindet:
\index{Funktionalgleichung}%
\[
\zeta(s) = 2^s \pi^{s-1} \sin\left(\frac{\pi s}{2}\right) \Gamma(1 - s)\,\zeta(1 - s).
\]
Die Gamma-Funktion $\Gamma$, die in dieser Formel auftritt, wird in Kapitel \ref{chapter:gamma} eingeführt.
\index{Gamma-Funktion@$\Gamma$-Funktion}%
Die Funktionalgleichung gilt für (fast) alle $s \in \mathbb{C}$ und liefert damit auch für negative $s$ einen wohldefinierten Wert von $\zeta(s)$, obwohl die ursprüngliche Reihe dort gar nicht konvergiert.

Setzt man $s = -1$ ein, so verbindet die Funktionalgleichung $\zeta(-1)$ mit $\zeta(2)$:
\[
\zeta(-1) = 2^{-1} \pi^{-2} \sin\left(-\frac{\pi}{2}\right)\Gamma(2)\,\zeta(2).
\]
Der Wert $\zeta(2)$ ist uns bereits bekannt: Es handelt sich, wie in Abschnitt~\ref{Riemannsche Zeta-Funktion} erwähnt, um das Basler Problem, welches Euler 1735 löste, mit dem Resultat $\zeta(2) = \pi^2/6$.
\index{Basel-Problem}%
\index{Euler, Leonhard}%

Mit
\[
\sin\left(-\frac{\pi}{2}\right) = -1, \quad \Gamma(2) = 1! = 1, \quad \zeta(2) = \frac{\pi^2}{6}
\]
folgt
\[
\zeta(-1) = \frac{1}{2} \cdot \frac{1}{\pi^2} \cdot (-1) \cdot 1 \cdot \frac{\pi^2}{6} = -\frac{1}{12}.
\]
Damit ist gezeigt, dass sich der Wert $-\frac{1}{12}$ nicht aus der divergenten Reihe selbst ergibt, sondern aus der Funktionalgleichung, die die analytische Fortsetzung der Zeta-Funktion liefert.




